\documentclass[a4paper,12pt]{article}
\usepackage[T2A]{fontenc}
\usepackage[utf8]{inputenc}
\usepackage[russian]{babel}
\usepackage{amsmath,amssymb,mathtools}
\usepackage{icomma}
\usepackage[a4paper,left=23mm,right=23mm,top=22mm,bottom=23mm,headheight=23pt]{geometry}
\usepackage{microtype}
\usepackage{setspace}
\usepackage{enumitem}
\usepackage{booktabs,tabularx}
\usepackage{tikz}
\usetikzlibrary{calc,arrows.meta,positioning,shapes.geometric}
\usepackage{xcolor}
\usepackage{fancyhdr}
\usepackage{titlesec}
\usepackage{hyperref}
\definecolor{accent}{HTML}{008080}
\definecolor{darkaccent}{HTML}{005656}
\definecolor{textgray}{HTML}{333333}
\definecolor{muted}{HTML}{777777}
\hypersetup{colorlinks=true,linkcolor=darkaccent,pdfauthor={Лёвин Артём Александрович},pdftitle={Чек-лист: корни, степени и преобразование выражений},pdfsubject={Математика. Анна Трапезникова. 17 ноября 2026 г.}}
\setstretch{1.12}
\setlength{\parindent}{0pt}
\setlength{\parskip}{5pt}
\setlist[itemize]{leftmargin=1.5em,itemsep=2pt,topsep=3pt}
\setlist[enumerate]{leftmargin=1.65em,itemsep=4pt,topsep=4pt}
\titleformat{\section}{\large\bfseries\color{darkaccent}}{\thesection.}{0.55em}{}
\titlespacing*{\section}{0pt}{1.03em}{0.44em}
\titleformat{\subsection}{\normalsize\bfseries\color{darkaccent}}{\thesubsection.}{0.5em}{}
\titlespacing*{\subsection}{0pt}{0.8em}{0.25em}
\pagestyle{fancy}
\fancyhf{}
\fancyhead[L]{\footnotesize\color{muted} Корни, степени и преобразования}
\fancyhead[R]{\footnotesize\color{muted} Анна Трапезникова}
\fancyfoot[C]{\footnotesize\thepage}
\renewcommand{\headrulewidth}{0.3pt}
\renewcommand{\headrule}{\hbox to\headwidth{\color{accent}\leaders\hrule height \headrulewidth\hfill}}
\newcommand{\keyidea}[1]{\par\smallskip{\color{darkaccent}\bfseries Ключевая идея.}\ #1\par}
\newcommand{\attention}[1]{\par\smallskip{\color{darkaccent}\bfseries Проверьте.}\ #1\par}
\newcommand{\worked}[1]{\par\smallskip{\color{darkaccent}\bfseries Пример.}\ #1\par}
\begin{document}
\thispagestyle{empty}
\begin{titlepage}
\centering
\vspace*{29mm}
{\color{darkaccent}\large МАТЕМАТИКА\par}
\vspace{21mm}
{\LARGE Чек-лист\par}
\vspace{13mm}
{\Large\color{accent} Корни, степени\\[2mm] и преобразование выражений\par}
\vspace{19mm}
{\normalsize Подготовка к ЕГЭ по профильной математике\par}
\vfill
{\large Анна Трапезникова\par}
\vspace{9mm}
{17 ноября 2026 г.\par}
\vfill
{\small Лёвин Артём Александрович\\ эксклюзивно для Анны Трапезниковой\par}
\vspace{15mm}
\begin{tikzpicture}
\draw[color=accent,line width=0.85pt] (0,0) .. controls (1.3,0.45) and (2.35,-0.45) .. (3.8,0)
 .. controls (5.15,0.45) and (6.3,-0.4) .. (7.6,0);
\end{tikzpicture}
\end{titlepage}
\setcounter{page}{1}

\section{Главный маршрут: как работать с корнями}
\keyidea{Большое или дробное число часто удобно представить как степень. После этого применяются правила степеней и корней.}
\[
\boxed{\text{Число} \ \longrightarrow\ \text{степень} \ \longrightarrow\ \text{свойство корня} \ \longrightarrow\ \text{ответ}}
\]
\textbf{Запомните числовые представления:}
\[
64=2^6,\quad 243=3^5,\quad 729=3^6,\quad
\frac1{125}=5^{-3},\quad 0{,}125=\left(\frac12\right)^3.
\]
Если число кажется неудобным, сначала попробуйте разложить его на простые множители или перевести десятичную дробь в обыкновенную.

\section{Корень и дробный показатель}
Арифметический корень чётной степени по определению \textbf{неотрицателен}. Для $a\ge0$ и натурального $n\ge2$:
\[
\sqrt[n]{a}=a^{1/n},\qquad
\sqrt[n]{a^m}=a^{m/n}.
\]
Для корня нечётной степени подкоренное выражение может быть и отрицательным:
\[
\sqrt[3]{-8}=-2, \qquad \sqrt[5]{-32}=-2.
\]
\worked{Вычислим $\sqrt[6]{36^3}$:}
\[
\sqrt[6]{36^3}=\sqrt[6]{(6^2)^3}
=\sqrt[6]{6^6}=6.
\]
\worked{Дробное основание:}
\[
\sqrt[4]{\left(\frac1{25}\right)^2}
=\sqrt[4]{\left(\frac15\right)^4}=\frac15.
\]

\section{Свойства степеней: короткая памятка}
При допустимых значениях выражений:
\[
\begin{aligned}
a^m a^n&=a^{m+n}, &\qquad \frac{a^m}{a^n}&=a^{m-n}\quad (a\ne0),\\
(a^m)^n&=a^{mn}, &(ab)^n&=a^n b^n,\\
a^{-n}&=\frac1{a^n}\quad(a\ne0), &a^{1/n}&=\sqrt[n]{a}\quad(a\ge0).
\end{aligned}
\]
\attention{При \emph{умножении} степеней с одинаковыми основаниями показатели складывают. При возведении степени в степень показатели перемножают.}
\[
2^4\cdot2^5=2^9,\qquad (2^4)^5=2^{20}.
\]

\section{Знак под корнем и модуль}
Для любого действительного $a$:
\[
\sqrt[2k]{a^{2k}}=|a|,\qquad
\sqrt[2k+1]{a^{2k+1}}=a\qquad(k\in\mathbb N).
\]
Поэтому
\[
\sqrt{x^2}=|x|,\qquad \sqrt[4]{(-3)^4}=3,
\qquad\sqrt[3]{(-3)^3}=-3.
\]
\keyidea{Минус можно вынести за знак корня \emph{нечётной} степени:}
\[
\sqrt[3]{-125}=-\sqrt[3]{125}=-5.
\]
Корень чётной степени из отрицательного числа в действительных числах не определён: например, $\sqrt[4]{-16}$.

\section{Уравнения вида \texorpdfstring{$x^n=a$}{x в степени n равно a}}
Число решений зависит от чётности $n$ и знака $a$.
\begin{itemize}
\item $n$ нечётное: для любого $a$ имеется ровно один действительный корень.
\item $n$ чётное и $a>0$: два корня, противоположные по знаку.
\item $n$ чётное и $a=0$: один корень $x=0$.
\item $n$ чётное и $a<0$: действительных корней нет.
\end{itemize}
\worked{В уравнении $x^4=256$ получаем $256=4^4$, поэтому $x=-4$ или $x=4$.}
\worked{В уравнении $x^5=-\frac1{32}$ получаем $-\frac1{32}=\left(-\frac12\right)^5$, поэтому $x=-\frac12$.}
\attention{Фраза «при чётной степени всегда два корня» верна только для положительной правой части.}

\section{Уравнения с дробным показателем}
\keyidea{Чтобы выделить неизвестное, используйте показатель, обратный исходному. Сначала проверьте область определения и возможные потери решений.}
\worked{Решим $x^{5/3}=32$:}
\[
\begin{aligned}
x&=32^{3/5}=(2^5)^{3/5}=2^3=\boxed{8}.
\end{aligned}
\]
\worked{Решим $t^{3/5}=8$:}
\[
 t=8^{5/3}=(2^3)^{5/3}=2^5=\boxed{32}.
\]
Здесь знаменатели показателей нечётные, поэтому обе степени определены для всех действительных значений переменных. В обоих примерах обратное преобразование сохраняет единственный корень.

\textbf{Осторожно с чётными числителями:} из $x^{2/3}=4$ следуют \emph{два} действительных решения: $x=-8$ и $x=8$, поскольку
\[
\left(\sqrt[3]{x}\right)^2=4.
\]

\section{Корень произведения и корень частного}
Для корня чётной степени при $a\ge0$, $b\ge0$:
\[
\sqrt[n]{ab}=\sqrt[n]{a}\,\sqrt[n]{b},
\qquad
\sqrt[n]{\frac ab}=\frac{\sqrt[n]{a}}{\sqrt[n]{b}}\quad(b>0).
\]
Для нечётного $n$ аналогичные равенства справедливы при любых действительных $a,b$ (в частном $b\ne0$).

\worked{Вместо перемножения чисел под корнем вычислим корни отдельно:}
\[
\sqrt[3]{343\cdot0{,}125}
=\sqrt[3]{7^3}\,\sqrt[3]{(0{,}5)^3}
=7\cdot0{,}5=\boxed{3{,}5}.
\]
\keyidea{Обратный ход тоже работает: если отдельные корни неудобны, объедините их под одним знаком.}
\[
\sqrt[3]{25}\cdot\sqrt[3]{40}=\sqrt[3]{1000}=10.
\]
\attention{Не раскладывайте корень \emph{суммы}: в общем случае $\sqrt{a+b}\ne\sqrt a+\sqrt b$.}

\section{Корни из десятичных дробей}
Удобно представить десятичную дробь как обыкновенную, а знаменатель --- как степень числа 10.
\[
\sqrt[3]{-0{,}001}
=-\sqrt[3]{\frac1{1000}}
=-\sqrt[3]{\left(\frac1{10}\right)^3}
=\boxed{-0{,}1}.
\]
\[
\sqrt[4]{0{,}0001}
=\sqrt[4]{\frac1{10^4}}=\boxed{0{,}1}.
\]

\section{Вложенные корни}
Для $a\ge0$ и натуральных $m,n\ge2$:
\[
\boxed{\sqrt[m]{\sqrt[n]{a}}=\sqrt[mn]{a}}.
\]
\worked{Вычислим вложенный корень:}
\[
\sqrt{\sqrt[3]{729}}
=\sqrt[6]{729}=\sqrt[6]{3^6}=\boxed{3}.
\]
\attention{Если под корнями есть переменные, после преобразования корня чётной степени проверьте, нужен ли модуль.}
\[
\sqrt{\sqrt[3]{x^6}}=\sqrt[6]{x^6}=|x|.
\]

\section{Буквенные выражения и ограничения}
Буквенные выражения преобразуют теми же свойствами, но при чётном корне важно сохранять неотрицательное значение результата.
\worked{При $b>0$ и любом действительном $a$:}
\[
\frac{\sqrt[4]{16a^4b^4}}{2b}
=\frac{2|a|\,|b|}{2b}
=\frac{2|a|b}{2b}
=\boxed{|a|}.
\]
Если известно дополнительно, что $a\ge0$, ответ можно записать как $a$.

\section{Частые ошибки и самопроверка}
\begin{itemize}
\item \textbf{Забыли модуль:} $\sqrt{x^2}=|x|$.
\item \textbf{Потеряли знак:} $\sqrt[3]{-27}=-3$.
\item \textbf{Пропустили второй корень:} $x^2=9\Rightarrow x=\pm3$.
\item \textbf{Перепутали действия с показателями:} $a^m a^n=a^{m+n}$, но $(a^m)^n=a^{mn}$.
\item \textbf{Применили свойство корня без условий:} для чётного корня проверьте знаки подкоренных выражений.
\item \textbf{Сократили переменную в знаменателе:} сначала убедитесь, что знаменатель не равен нулю.
\end{itemize}
\keyidea{Сначала установите допустимые значения и удобную форму числа; в конце проверьте знак результата и возможное появление модуля.}

\clearpage
\thispagestyle{fancy}
\begin{center}
{\Large\color{darkaccent}\bfseries Алгоритм: преобразование корней и степеней}
\end{center}
\vspace{15mm}
\begin{center}
\begin{tikzpicture}[
 >={Latex[length=2.6mm]},
 endnode/.style={ellipse,draw=accent,line width=.85pt,align=center,text width=7.0cm,minimum height=1.05cm,fill=white},
 action/.style={rounded corners=2pt,draw=accent,line width=.7pt,align=center,text width=5.45cm,minimum height=1.45cm,inner sep=6pt,fill=white},
 decision/.style={diamond,draw=darkaccent,line width=.85pt,aspect=2.5,align=center,text width=4.6cm,inner xsep=1pt,inner ysep=3pt,fill=white},
 connector/.style={-Latex,line width=.8pt,color=darkaccent}
]
\node[endnode] (start) at (0,0) {Дано выражение с корнями и степенями};
\node[decision] (dec) at (0,-2.3) {Нужно решить уравнение?};
\node[action] (left1) at (-3.92,-5.2) {Представьте составные числа и дроби как степени. Установите условия существования корней.};
\node[action] (right1) at (3.92,-5.2) {Определите чётность степени, знак правой части и область допустимых значений.};
\node[action] (left2) at (-3.92,-8.0) {Разделите или объедините корни, упростите вложенные корни и показатели.};
\node[action] (right2) at (3.92,-8.0) {Выделите неизвестную степень. При дробном показателе проверьте число решений.};
\node[endnode] (end) at (0,-10.9) {Проверьте знаки, модули, ограничения и запишите итог};
\draw[connector] (start.south)--(dec.north);
\draw[connector] (dec.west)-| node[pos=.27,above,font=\small] {нет} (left1.north);
\draw[connector] (dec.east)-| node[pos=.27,above,font=\small] {да} (right1.north);
\draw[connector] (left1.south)--(left2.north);
\draw[connector] (right1.south)--(right2.north);
\draw[connector] (left2.south) |- ([yshift=6mm]end.west) -- (end.west);
\draw[connector] (right2.south) |- ([yshift=6mm]end.east) -- (end.east);
\end{tikzpicture}
\end{center}

\clearpage
\section*{Тренировочный блок}
\textbf{Путь А. 10 заданий --- от базовых преобразований к комбинированным.}

Вычислите, решите уравнения или упростите выражения. В заданиях с буквами внимательно следите за ограничениями и модулями. Решения записывайте последовательно.
\vspace{3mm}
\begin{enumerate}[label=\textbf{\arabic*.},leftmargin=2.2em,itemsep=0.67\baselineskip]
\item Вычислите:
\[ \sqrt[3]{125}. \]
\item Вычислите:
\[ \sqrt[4]{16^3}. \]
\item Вычислите:
\[ \sqrt[4]{\left(\frac1{25}\right)^2}. \]
\item Вычислите:
\[ \sqrt[3]{-0{,}027}. \]
\item Решите уравнение:
\[ x^4=1296. \]
\item Решите уравнение:
\[ x^5=-\frac1{32}. \]
\item Решите уравнение:
\[ x^{3/5}=27. \]
\item Вычислите, объединив корни:
\[ \sqrt[3]{40}\cdot\sqrt[3]{25}. \]
\item Упростите при $b>0$ ($a$ --- любое действительное число):
\[ \frac{\sqrt[4]{16a^4b^4}}{2b}. \]
\item Упростите при любом действительном $x$:
\[ \sqrt{\sqrt[3]{729x^6}}. \]
\end{enumerate}

\clearpage
\section*{Ответы к тренировочному блоку}
\vspace{4mm}
\renewcommand{\arraystretch}{1.58}
\begin{center}
\begin{tabularx}{0.87\linewidth}{@{}>{\centering\arraybackslash}p{0.17\linewidth}X@{}}
\toprule
\textbf{№} & \textbf{Ответ}\\
\midrule
1 & $5$\\
2 & $8$\\
3 & $\dfrac15$\\
4 & $-0{,}3$\\
5 & $x=-6$ или $x=6$\\
6 & $x=-\dfrac12$\\
7 & $x=243$\\
8 & $10$\\
9 & $|a|$\\
10 & $3|x|$\\
\bottomrule
\end{tabularx}
\end{center}
\vspace{9mm}
\section*{Итоговая памятка}
\begin{tabularx}{\linewidth}{@{}p{0.34\linewidth}X@{}}
\toprule
\textbf{Что Вы видите} & \textbf{Что полезно сделать}\\
\midrule
Составное число & Представить степенью простого числа.\\
Десятичная дробь & Перевести в обыкновенную дробь.\\
Произведение под корнем & Разделить на корни, если выполнены условия.\\
Произведение корней & Объединить под одним корнем.\\
Корень под корнем & Перемножить показатели корней.\\
Чётная степень в уравнении & Проверить знак правой части и оба корня.\\
Корень чётной степени из $a^{2k}$ & Не потерять $|a|$.\\
Дробный показатель & Использовать обратную степень с учётом области определения и числа решений.\\
\bottomrule
\end{tabularx}
\vfill
\begin{center}
{\small\color{muted}Лёвин Артём Александрович --- эксклюзивно для Анны Трапезниковой}
\end{center}
\end{document}
